\begin{ouverture}[geometrie]{12}{Pythagore, cercle et droite des milieux}{Partie 4 · Espace et géométrie}
{Pythagore · Réciproque · Contraposée · Cercle circonscrit · Droite des milieux}
\objectif{Calculer une longueur dans un triangle rectangle.}
\objectif{Prouver ou réfuter qu’un triangle est rectangle.}
\objectif{Utiliser le cercle circonscrit et les théorèmes des milieux.}
\prerequis{Carrés, racines carrées, médiatrices et parallélogrammes.}
\end{ouverture}
\pagestyle{otMain}\essentiel\fontsize{11}{15}\selectfont
\essSection{1}{Pythagore, réciproque et contraposée}
\begin{propriete}[Théorème de Pythagore]Si ABC est rectangle en A, alors $BC^2=AB^2+AC^2$. BC, opposé à l’angle droit, est l’hypoténuse et le plus long côté.\end{propriete}
\begin{savoirfaire}[Calculer une longueur]Pour $AB=6$ cm et $AC=8$ cm, $BC^2=36+64=100$, donc $BC=\sqrt{100}=10$ cm. Si $BC=13$ cm et $AB=5$ cm, $AC^2=169-25=144$, d’où $AC=12$ cm. Une longueur est positive.\end{savoirfaire}
\begin{propriete}[Réciproque]Dans un triangle, si le carré du plus long côté est égal à la somme des carrés des deux autres, alors le triangle est rectangle, et l’angle droit est opposé au plus long côté.\end{propriete}
\begin{definition}La réciproque de « si $P$, alors $Q$ » est « si $Q$, alors $P$ ». Sa contraposée est « si $Q$ est faux, alors $P$ est faux ». Une propriété vraie a une contraposée vraie ; sa réciproque doit être démontrée à part.\end{definition}
\begin{propriete}[Contraposée de Pythagore]Si le carré du plus long côté n’est pas égal à la somme des carrés des deux autres, alors le triangle n’est pas rectangle.\end{propriete}
\begin{attention}Le théorème exige un angle droit connu : une apparence de dessin ne suffit pas. Conserver la valeur exacte d’une racine, puis arrondir à la précision demandée.\end{attention}
\clearpage
\essSection{2}{Cercle circonscrit et droite des milieux}
\begin{propriete}Si ABC est rectangle en A, alors A appartient au cercle de diamètre [BC] : le centre du cercle circonscrit est le milieu de l’hypoténuse. Réciproquement, si A appartient au cercle de diamètre [BC], avec A distinct de B et de C, alors ABC est rectangle en A.\end{propriete}
\begin{propriete}[Droite des milieux]Dans un triangle ABC :
\par 1. la droite qui passe par les milieux de [AB] et de [AC] est parallèle à (BC) ;
\par 2. la droite qui passe par le milieu de [AB] parallèlement à (BC) coupe [AC] en son milieu ;
\par 3. le segment qui joint les milieux de [AB] et de [AC] mesure la moitié de BC.\end{propriete}
\begin{exemple}Si I et J sont les milieux de [AB] et [AC] et si $BC=9$ cm, alors $(IJ)\parallel(BC)$ et $IJ=4{,}5$ cm.\end{exemple}
\begin{savoirfaire}[Rédiger]Citer les hypothèses (angle droit, milieux, points sur un cercle), nommer la propriété utilisée, puis conclure en nommant précisément les objets. Pour montrer qu’un triangle n’est pas rectangle, le mot « contraposée » est attendu.\end{savoirfaire}
\clearpage\rubrique{Découverte}
\begin{blocLivre}{ex-12-1}\begin{exercice}[1]{Reconnaître l’hypoténuse}
ABC est rectangle en A. Nommer l’hypoténuse et écrire l’égalité de Pythagore.
\label{fig-ex-12-1}\figureGeo{pyth-01}
\end{exercice}\end{blocLivre}
\begin{blocLivre}{ex-12-2}\begin{exercice}[1]{Un triangle classique}
ABC est rectangle en A, avec AB = 6 cm et AC = 8 cm. Calculer BC.
\label{fig-ex-12-2}\figureGeo{pyth-02}
\end{exercice}\end{blocLivre}
\begin{blocLivre}{ex-12-3}\begin{exercice}[1]{Le côté manquant}
ABC est rectangle en A. BC = 13 cm et AB = 5 cm. Calculer AC.
\label{fig-ex-12-3}\figureGeo{pyth-03}
\end{exercice}\end{blocLivre}
\begin{blocLivre}{ex-12-4}\begin{exercice}[1]{Garder une racine}
ABC est rectangle en A. AB = 4 cm, AC = 7 cm. Donner BC exactement puis au dixième.
\label{fig-ex-12-4}\figureGeo{pyth-04}
\end{exercice}\end{blocLivre}
\begin{blocLivre}{ex-12-5}\begin{exercice}[1]{Le centre du cercle circonscrit}
ABC est rectangle en A, avec BC = 10 cm ; O est le milieu de [BC]. Où se trouve le centre du cercle circonscrit à ABC ? Donner son rayon et la longueur OA.
\label{fig-ex-12-5}\figureGeo{cercle-circonscrit}
\end{exercice}\end{blocLivre}

\clearpage\rubrique{Entraînement}
\begin{blocLivre}{ex-12-6}\begin{exercice}[2]{Tester la réciproque}
Dans ABC, AB = 9 cm, AC = 12 cm et BC = 15 cm. Le triangle est-il rectangle ?
\label{fig-ex-12-6}\figureGeo{pyth-05}
\end{exercice}\end{blocLivre}
\begin{blocLivre}{ex-12-7}\begin{exercice}[2]{Utiliser la contraposée}
AB = 5 cm, AC = 7 cm et BC = 9 cm. ABC est-il rectangle ? Nommer la propriété utilisée.
\label{fig-ex-12-7}\figureGeo{pyth-06}
\end{exercice}\end{blocLivre}
\begin{blocLivre}{ex-12-8}\begin{exercice}[2]{La droite des milieux}
Dans le triangle ABC, I est le milieu de [AB] et J celui de [AC] ; BC = 9 cm. Que peut-on dire des droites (IJ) et (BC) ? Calculer IJ.
\label{fig-ex-12-8}\figureGeo{milieux-parallele}
\end{exercice}\end{blocLivre}
\begin{blocLivre}{ex-12-9}\begin{exercice}[2]{Un milieu retrouvé}
Dans le triangle ABC, I est le milieu de [AB]. La parallèle à (BC) passant par I coupe [AC] en J. On sait que AC = 7,4 cm. Calculer AJ en justifiant.
\label{fig-ex-12-9}\figureGeo{milieux-milieu}
\end{exercice}\end{blocLivre}
\begin{blocLivre}{ex-12-10}\begin{exercice}[2]{Un angle à identifier}
AB = 10 cm, AC = 6 cm et BC = 8 cm. Prouver que le triangle est rectangle et préciser en quel sommet.
\label{fig-ex-12-10}\figureGeo{pyth-10}
\end{exercice}\end{blocLivre}

\clearpage\rubrique{Problèmes}
\begin{blocLivre}{ex-12-11}\begin{exercice}[3]{Une échelle contre un mur}
Le mur AC est vertical et le sol AB horizontal. Une échelle BC de 5 m est posée à 1,4 m du mur. Quelle hauteur AC atteint-elle ?
\label{fig-ex-12-11}\figureGeo{pyth-07}
\end{exercice}\end{blocLivre}
\begin{blocLivre}{ex-12-12}\begin{exercice}[3]{Un triangle dans un demi-cercle}
[BC] est un diamètre d’un cercle de centre O et de rayon 6,5 cm. A est un point de ce cercle, distinct de B et de C, tel que AB = 5 cm. Démontrer que ABC est rectangle, puis calculer AC.
\label{fig-ex-12-12}\figureGeo{demi-cercle}
\end{exercice}\end{blocLivre}
\begin{blocLivre}{ex-12-13}\begin{exercice}[3]{Un rectangle sans équerre}
On trace un cercle de centre O, puis deux diamètres [AC] et [BD] qui ne sont pas perpendiculaires. Démontrer que ABCD est un rectangle. Pourquoi ce procédé permet-il de construire un rectangle sans équerre ?
\label{fig-ex-12-13}\figureGeo{rectangle-compas}
\end{exercice}\end{blocLivre}
\begin{blocLivre}{ex-12-14}\begin{exercice}[3]{Une hauteur de triangle}
ABC est isocèle en C : CA = CB = 13 cm et AB = 10 cm. H est le milieu de [AB]. Calculer CH puis l’aire de ABC.
\label{fig-ex-12-14}\figureGeo{pyth-iso}
\end{exercice}\end{blocLivre}
\begin{blocLivre}{ex-12-15}\begin{exercice}[3]{Le théorème de Varignon}
ABCD est un quadrilatère quelconque, non croisé. I, J, K et L sont les milieux respectifs de [AB], [BC], [CD] et [DA]. Démontrer que IJKL est un parallélogramme.
\label{fig-ex-12-15}\figureGeo{varignon}
\end{exercice}\end{blocLivre}

\clearpage
